\section{Construction: enlarge the closed region, then share a selector}
\label{sec:construction}

The earlier tagged-selector proposal could not simply be stacked with U212:
it repurposed native helper 211, which the accepted U212 patch enters from
$180/1$. The successful construction preserves that helper and changes a
larger closed source region. No axiom is removed and no new proof-search
calculus is introduced.

\subsection{Source reassociation and physical decoder rows}

The source is a full binary sequence tree. Leaves are register primitives
with effects \code{CI}, \code{CD}, \code{CB}, \code{CC}, or \code{CH}; a leaf's
binary address identifies its position in the tree. Branch primitives
refer to enclosing left or right ancestors, so arbitrary reassociation
would not be justified. A factor is \emph{closed} when it has no branch
escaping its root and contains no halt primitive. For three closed factors,
the changes
\[
        ((A,B),C)\quad\longleftrightarrow\quad(A,(B,C))
\]
preserve the ordered primitive effects and their resolved ordinal successors.
The checker establishes closure for the exact factors and then verifies
the entire ordinal control-flow graph, including both cases of every
conditional decrement.

A pure physical decoder state reads a bit, writes that same bit, moves
right, and enters its zero or one child. Write its children as an ordered
pair. If $p=(q,c)$ and $q=(a,b)$, a left rotation makes
$p=(a,q)$ and $q=(b,c)$. If $p=(a,q)$ and $q=(b,c)$, the inverse rotation
makes $p=(q,c)$ and $q=(a,b)$. In both cases the inner decoder $q$ must
have precisely the stated incoming row. This is checked over the whole
literal table, not merely the source tree. The builder also rejects any
incoming left-moving runtime edge to either edited decoder.

The source addresses are changed accordingly. For the first direction,
at each affected prefix $s$, the replacements are
\[
       s00u\mapsto s0u,\qquad
       s01u\mapsto s10u,\qquad
       s1u\mapsto s11u;
\]
the reverse direction uses the inverse prefix code. All occurrences of a
shared physical decoder are transformed. The independent auditor reconstructs
these occurrences from literal prefix executions rather than trusting
the discovery program's map. At every stage it checks all 3,908 internal
prefixes, rejecting a write that changes a bit or a read outside the
known prefix and left sentinel. It independently rechecks the full source
graph after the transformation.

\subsection{The nine rotations}

The useful closed region is rooted at native 174 and occupies primitive
ordinals $[2722,3658)$, with zero-based ordinals. It extends the previously
used 774-primitive emitter by one preceding closed 162-primitive factor.
Its fourteen factors, in unchanged order, are
\[
 159,147,176,98,218,168,98,147,146,218,176,90,211,218.
\]
Their primitive lengths are respectively
\[
 162,52,80,26,56,52,26,52,108,56,80,27,103,56,
\]
which sum to 936. Repeated labels denote repeated occurrences of the same
factor. The nine literal rotations are listed in Table~\ref{tab:rotations}.

\begin{table}[htbp]
\centering
\begin{tabular}{rrrrl}
\toprule
Step & Outer & Inner & Side & Source reassociation\\
\midrule
1&215&204&1&$(A,(B,C))\to((A,B),C)$\\
2&217&181&1&$(A,(B,C))\to((A,B),C)$\\
3&185&217&0&$((A,B),C)\to(A,(B,C))$\\
4&174&185&1&$(A,(B,C))\to((A,B),C)$\\
5&129&158&0&$((A,B),C)\to(A,(B,C))$\\
6&217&215&1&$(A,(B,C))\to((A,B),C)$\\
7&204&214&0&$((A,B),C)\to(A,(B,C))$\\
8&215&129&0&$((A,B),C)\to(A,(B,C))$\\
9&216&220&1&$(A,(B,C))\to((A,B),C)$\\
\bottomrule
\end{tabular}
\caption{The deterministic recipe. Side is the old edge from outer to inner.
All labels are native. Exact occurrence prefixes are in
\protect\path{inputs/rotation_paths.json}.}
\label{tab:rotations}
\end{table}

One way to display the final association, using $(x,y)$ for source
sequencing, is
\begin{align*}
 F&=(211,218), & E&=(218,176),\\
 D&=(146,(E,90)), & C&=(147,D),\\
 B&=((218,(168,98)),C), & A&=(98,B),\\
 \text{root }174&=((159,(147,176)),(A,F)).
\end{align*}
In the physical table, $B$ uses helper 129; its left branch through 158
begins with 218 and its right branch through 204 begins with 147. The
factor represented by 211 remains intact.

After all nine rotations, the source still has 3,909 primitives and 197
distinct subtrees, maximum depth 42 and maximum zero run six. The reduction
does not improve the subtree count: it makes the existing allocation more
shareable. Intermediate rotated tables are construction stages only. Some
intermediate address shapes are unsuitable for the original scanner; the
full runtime theorem is checked for the final table, not asserted for
every intermediate stage.

\section{The predecessor-tag fold}
\label{sec:fold}

Let $\Rot$ denote the 212-state table after the nine rotations. Its relevant
rows are:
\begin{center}
\begin{tabular}{rcc}
\toprule
Native state & Read 0 & Read 1\\
\midrule
220 & $\act{0}{R}{218}$ & $\act{1}{R}{176}$\\
181 & $\act{0}{R}{147}$ & $\act{1}{R}{176}$\\
129 & $\act{0}{R}{158}$ & $\act{1}{R}{204}$\\
158 & $\act{0}{R}{218}$ & unchanged\\
204 & $\act{0}{R}{147}$ & unchanged\\
\bottomrule
\end{tabular}
\end{center}
The complete incoming inventories, including all runtime rows, are
\[
 \operatorname{In}(220)=\{216/0:\act{0}{R}{220}\},\qquad
 \operatorname{In}(181)=\{185/1:\act{1}{R}{181}\}.
\]
Neither state is initial. Thus whenever $\Rot$ enters 220, the cell immediately
left of its head is 0; whenever it enters 181, that cell is 1. The previous
transition wrote the identifying bit, so no condition on earlier tape
contents is needed.

The final fold performs exactly three operations:
\begin{enumerate}
\item Redirect the target of $185/1$ from 181 to 220.
\item Replace $220/0$ by $\act{0}{L}{129}$.
\item Delete the two outgoing rows of 181.
\end{enumerate}
All other retained rows agree with $\Rot$, up to the target identification
$\rho(181)=220$, $\rho(q)=q$ otherwise. In particular, the helpers are checked
in the final edited table, and are distinct from both merged states.

\begin{lemma}[Four local identities]\label{lem:four}
Let the old head position be $h$, with left-neighbor tag 0 at old state 220
or tag 1 at old state 181. One old transition is simulated by one or three
new transitions, with the same whole tape and head position and control
mapped by $\rho$. The exterior tape is arbitrary and is not read.
\end{lemma}
\begin{proof}
The two possible current symbols give the four cases in
Table~\ref{tab:local}. If the current symbol is 1, both old states have
the identical row to 176, which is retained. If it is 0, the new 220 row
moves left to 129. That helper reads and preserves the tag, returning right
to 158 or 204. The current zero then selects the required old target 218
or 147. Every write preserves the symbol read. Only cells $h-1$ and $h$
are read; the head ends at $h+1$, which is not read inside the block.
No displayed intermediate state is halt.
\end{proof}

\begin{table}[htbp]
\centering\small
\begin{tabular}{rrrlc}
\toprule
Old state & Tag & Read & New state/head path, relative to $h$ & Steps\\
\midrule
220&0&0&$220@0\to129@{-1}\to158@0\to218@1$&3\\
181&1&0&$220@0\to129@{-1}\to204@0\to147@1$&3\\
220&0&1&$220@0\to176@1$&1\\
181&1&1&$220@0\to176@1$&1\\
\bottomrule
\end{tabular}
\caption{Exhaustive local cases. The old block has one step in every row.}
\label{tab:local}
\end{table}

\begin{theorem}[Final-fold equivalence]\label{thm:fold}
For any common initial tape, from start state 0 and a common head position,
$\Rot$ halts if and only if $\New$ halts. If a common halting run has
$T_{\mathrm{rot}}$ old steps and $k$ of them read zero at old 220 or 181,
then
\[
 T_{211}=T_{\mathrm{rot}}+2k,\qquad
 T_{\mathrm{rot}}\leq T_{211}\leq3T_{\mathrm{rot}}.
\]
\end{theorem}
\begin{proof}
At time zero neither tagged state is active. The complete incoming lists
establish the required tag at every subsequent old occurrence. For every
old step outside the two tagged states, the corresponding new row commutes
with $\rho$. There are 420 such rows, all checked explicitly. The common
one branch and Lemma~\ref{lem:four} handle the remaining rows.

Induction on old time therefore embeds the old execution into the unique
new execution at block boundaries, with identical tape and head and
projected control. Every block is finite, has length one or three, and
has no internal halt. If the old run halts, the corresponding new boundary
halts. If the old run is infinite, its block times grow without bound and
describe a nonhalting new run; a new halt cannot occur inside a block.
This proves reflection as well as preservation. Summing block lengths
gives the exact time identity and bounds.
\end{proof}

The arbitrary-tape statement is an ordinary mathematical local simulation.
For its PA formulation one codes finite configurations and finite tape
support; blank-tape equivalence and the finite trace induction are arithmetic.
No second-order quantification over infinite tapes is required. The bound
belongs only to the final fold $\Rot\to\New$. The preceding reassociations
change source addresses, so neither this bound nor arbitrary-common-tape
equivalence is asserted for the whole path from the original $\Old$.

The finite structural verifier \path{code/verify_fold.py} checks the
incoming inventories, all ordinary rows, the final helper rows, and the
four paths over partial tapes with an unread opaque exterior. Twelve
intentional corruptions are rejected. The verifier uses explicit checks
and passes also under Python \code{-O}; it does not turn off the obligations
when assertions are disabled.
